\documentclass{article}
\usepackage[utf8]{inputenc}
\usepackage{amsmath}
\usepackage{amssymb}
\usepackage{booktabs}
\usepackage{listings}
\usepackage{xcolor}

\definecolor{codegreen}{rgb}{0,0.6,0}
\definecolor{codegray}{rgb}{0.5,0.5,0.5}
\definecolor{codepurple}{rgb}{0.58,0,0.82}
\definecolor{backcolour}{rgb}{0.95,0.95,0.92}

\lstset{
    backgroundcolor=\color{backcolour},   
    commentstyle=\color{codegreen},
    keywordstyle=\color{magenta},
    numberstyle=\tiny\color{codegray},
    stringstyle=\color{codepurple},
    basicstyle=\ttfamily\footnotesize,
    breaklines=true,                 
    captionpos=b,                    
    keepspaces=true,                 
    numbers=left,                    
    numbersep=5pt,                  
    showspaces=false,                
    showstringspaces=false,
    showtabs=false,                  
    tabsize=2
}

\title{Updated 6DoF Sensing Framework: Scaling Limits and Spine Density}
\author{Analysis of 3mm Miniaturization Tax and Continuum Reconstruction}
\date{}

\begin{document}

\maketitle

\section{The "Miniaturization Tax" and Noise Amplification}
Experimental modeling of the 3mm sensor board reveals a critical stability threshold. As the board diameter $D$ decreases, the condition number $\kappa(\mathbf{M})$ grows exponentially. 

\begin{itemize}
    \item \textbf{Rotation Error ($R_{err}$):} Below 5mm, the leverage arm $L$ is insufficient to distinguish tilt from translation. At 3mm, rotational noise is amplified by approximately $4\times$ compared to a 20mm baseline.
    \item \textbf{Translation Jitter:} While $X, Y, Z$ are less dependent on $L$, they suffer from high-frequency "jagged" noise due to the low Signal-to-Noise Ratio (SNR) of small-aperture photodiodes.
\end{itemize}

\section{Continuum Spine Reconstruction Density}
To reconstruct the shape of a soft robot (e.g., "pool noodle" analogy) bending up to $\theta = 75^\circ$ ($1.31$ rad) over a length $L_{spine} = 100$mm, we calculate the minimum bend radius $\rho_{min}$:

\begin{equation}
\rho_{min} = \frac{L_{spine}}{\theta_{rad}} \approx 76.3\text{ mm}
\end{equation}

Applying the Nyquist-criterion for curvature sampling, the sensor spacing $S$ must satisfy:
\begin{equation}
S \leq \frac{\rho_{min}}{2} \approx 38\text{ mm}
\end{equation}

\textbf{Recommended Configuration:} To account for S-curves and complex strain localized at the 3mm board, a spacing of \textbf{20--25mm} is recommended, resulting in a 5-sensor "string of pearls" for a 100mm spine.

\section{Mechanical Interaction: Rigid-Inclusion Correction}
Because the 3mm PCB is rigid relative to the 5mm soft spine, the sensor creates a "dead zone" where the silicone cannot bend naturally. The effective sensing length $L_{eff}$ is reduced by the inclusion diameter $D$:

\begin{equation}
L_{eff} = S - (1.5 \times D_{sensor})
\end{equation}

For a 3mm board, each sensor "blindfolds" approximately 4.5mm of the robot's local curvature.

\begin{lstlisting}[language=Matlab, caption=Optimal Density Calculation]
% Calculate required sensors for a given max bend
L_spine = 100; % mm
max_bend_deg = 75;
rho_min = L_spine / deg2rad(max_bend_deg);
spacing = rho_min / 2; % Nyquist limit
num_sensors = ceil(L_spine / spacing);
fprintf('Required Sensors: %d at %0.1f mm intervals\n', num_sensors, spacing);
\end{lstlisting}

\end{document}

